\documentclass[10pt,a4paper]{article}

\usepackage[utf8]{inputenc}
\usepackage[T1]{fontenc}
\usepackage[margin=2.5cm]{geometry}
\usepackage{hyperref}
\usepackage{booktabs}
\usepackage{array}
\usepackage{multirow}
\usepackage{graphicx}
\usepackage{caption}
\usepackage{subcaption}
\usepackage{amsmath}
\usepackage{enumitem}
\usepackage{xcolor}
\usepackage{parskip}

\hypersetup{
    colorlinks=true,
    linkcolor=blue,
    urlcolor=blue,
    citecolor=blue
}

\title{
    \textbf{Mid-Phase Progress Report} \\[0.5em]
    \large Comparative Study of Prompting Strategies for Secure Code Generation \\
    Using Local Large Language Models
}

\author{
    Ali Ekin Özçetin \\
    \texttt{ozcetinaliekin@gmail.com} \\[0.5em]
    \small CENG 467 — Natural Language Processing \\
    \small Forked from: \href{https://github.com/Aliekinozcetin/CodeEnhancer}{github.com/Aliekinozcetin/CodeEnhancer}
}

\date{May 2026}

\begin{document}

\maketitle

% ─────────────────────────────────────────────────────────────────────────────
\section{Project Title, Group Members, and Problem Definition}
% ─────────────────────────────────────────────────────────────────────────────

\textbf{Project Title:} Comparative Evaluation of Prompting Strategies for Reducing Security Vulnerabilities in LLM-Generated Python Code

\textbf{Group Member:} Ali Ekin Özçetin

\subsection*{Problem Definition}

LLM-generated code frequently contains security vulnerabilities \cite{pearce2022asleep,tony2023llmseceval}. This project investigates whether \emph{prompting strategy} (zero-shot, few-shot, chain-of-thought) and \emph{model choice} (Llama 3.1 8B, DeepSeek-Coder 6.7B, Mistral 7B) influence the frequency and resolvability of security issues in LLM-generated Python code.

The pipeline extends the CodeEnhancer Stage 1 framework \cite{nishida2024codeenhancer}: each generated file is subjected to Bandit + Pylint SAST, an LLM-driven repair loop (up to 5 attempts), and a final functional correctness check. Key metrics: Bandit hit rate at iteration 0, resolution rate, unresolved rate, average iterations to resolve, and CWE breakdown.

% ─────────────────────────────────────────────────────────────────────────────
\section{Dataset / Benchmark Status}
% ─────────────────────────────────────────────────────────────────────────────

The dataset used is \textbf{LLMSecEval} \cite{tony2023llmseceval}, a benchmark of 150 natural language prompts designed to elicit security-relevant Python code across 18 CWE categories drawn from the CWE Top 25.

\subsection*{Mid-Phase Subset}

15 prompts were selected via stratified sampling (seed=42) from the 83 Python-language prompts, ensuring at least one prompt per CWE category (\texttt{data/select\_prompts.py}).

\begin{table}[h]
\centering
\caption{CWE distribution of the 15 mid-phase prompts}
\begin{tabular}{lll}
\toprule
\textbf{CWE ID} & \textbf{Name} & \textbf{Count} \\
\midrule
CWE-78  & OS Command Injection        & 2 \\
CWE-89  & SQL Injection               & 2 \\
CWE-502 & Deserialization             & 2 \\
CWE-22  & Path Traversal              & 1 \\
CWE-79  & Cross-Site Scripting        & 1 \\
CWE-434 & Unrestricted File Upload    & 1 \\
CWE-306 & Missing Auth                & 1 \\
CWE-200 & Info Exposure               & 1 \\
CWE-20  & Improper Input Validation   & 1 \\
CWE-732 & Incorrect Permissions       & 1 \\
CWE-522 & Weak Credentials            & 1 \\
CWE-798 & Hardcoded Credentials       & 1 \\
\bottomrule
\end{tabular}
\label{tab:cwe_dist}
\end{table}

The selected prompts are saved in \texttt{data/mid\_phase\_prompts.json} and shared as a control variable across any parallel branches.

% ─────────────────────────────────────────────────────────────────────────────
\section{Literature Review Progress}
% ─────────────────────────────────────────────────────────────────────────────

Five papers directly inform the methodology: Pearce et al. \cite{pearce2022asleep} show GitHub Copilot produces vulnerable code in 40\% of scenarios, motivating iterative repair pipelines. Tony et al. \cite{tony2023llmseceval} introduce LLMSecEval and report a 72\% Bandit hit rate on ChatGPT, which is the direct baseline for comparison. Nishida et al. \cite{nishida2024codeenhancer} describe CodeEnhancer Stage 1, the architectural basis of this project. Brown et al. \cite{brown2020gpt3} justify using 3 few-shot examples as the standard, and Wei et al. \cite{wei2022cot} motivate the 4-step chain-of-thought structure.

Planned: Kojima et al. zero-shot CoT (2022) and HumanEval-Security for the final report.

% ─────────────────────────────────────────────────────────────────────────────
\section{Baseline Models and Prompting Strategies}
% ─────────────────────────────────────────────────────────────────────────────

All models run locally via Ollama on Apple M4 (16 GB, Metal GPU): \textbf{Llama 3.1 8B} (Meta, general-purpose), \textbf{DeepSeek-Coder 6.7B} (DeepSeek, code-specialized), \textbf{Mistral 7B} (Mistral AI, general-purpose). This covers one code-domain model vs.\ two general-purpose models.

Three strategies are implemented in \texttt{prompts/}: \textbf{Zero-shot} (original CodeEnhancer system prompt, no security instruction), \textbf{Few-shot} (3 secure examples covering CWE-78, CWE-89, CWE-502, following Brown et al.\ \cite{brown2020gpt3}), and \textbf{Chain-of-Thought} (4-step security reasoning template, following Wei et al.\ \cite{wei2022cot}).

This 3×3 design yields \textbf{9 conditions} × 15 prompts = \textbf{135 generated files}.

% ─────────────────────────────────────────────────────────────────────────────
\section{Initial Experimental Results}
% ─────────────────────────────────────────────────────────────────────────────

All 135 files were generated and validated through the full pipeline (SAST → repair loop, MAX\_ATTEMPTS=5 → functional check).

\subsection*{Per-Model Summary}

\begin{table}[h]
\centering
\caption{Aggregated metrics by model (averaged across 3 strategies, 15 prompts each)}
\begin{tabular}{lcccc}
\toprule
\textbf{Model} & \textbf{Bandit@0 (\%)} & \textbf{Resolution (\%)} & \textbf{Unresolved (\%)} & \textbf{Avg Iter} \\
\midrule
Llama 3.1 8B       & 44.4 & 15.6 & 84.4 & 4.73 \\
DeepSeek-Coder 6.7B & 28.9 & 20.0 & 80.0 & 4.53 \\
Mistral 7B         & 31.1 & 75.6 & 24.4 & 2.62 \\
\bottomrule
\end{tabular}
\label{tab:model_summary}
\end{table}

Mistral 7B substantially outperforms the others in resolution rate (75.6\% vs.\ 15.6\% / 20.0\%) and converges faster (2.62 vs.\ 4.73 avg iterations), suggesting stronger instruction-following in the repair loop.

\subsection*{Per-Strategy Summary}

\begin{table}[h]
\centering
\caption{Aggregated metrics by prompting strategy (averaged across 3 models, 15 prompts each)}
\begin{tabular}{lcccc}
\toprule
\textbf{Strategy} & \textbf{Bandit@0 (\%)} & \textbf{Resolution (\%)} & \textbf{Unresolved (\%)} & \textbf{Avg Iter} \\
\midrule
Zero-shot         & 42.2 & 35.5 & 64.5 & 4.09 \\
Few-shot          & 28.9 & 37.8 & 62.2 & 3.84 \\
Chain-of-Thought  & 33.3 & 37.8 & 62.2 & 3.96 \\
\bottomrule
\end{tabular}
\label{tab:strategy_summary}
\end{table}

Few-shot and CoT both reduce the initial Bandit hit rate vs.\ zero-shot. Resolution rates are similar across strategies (35.5–37.8\%), suggesting repair-loop performance is dominated more by model capability than prompting strategy.

\subsection*{Full Results Grid}

\begin{table}[h]
\centering
\caption{Complete results for all 9 experiment conditions}
\begin{tabular}{llcccc}
\toprule
\textbf{Model} & \textbf{Strategy} & \textbf{Bandit@0 (\%)} & \textbf{Resolution (\%)} & \textbf{Unresolved (\%)} & \textbf{Avg Iter} \\
\midrule
Llama 3.1 8B        & Zero-shot  & 66.7 & 13.3 & 86.7 & 4.73 \\
Llama 3.1 8B        & Few-shot   & 33.3 & 13.3 & 86.7 & 4.73 \\
Llama 3.1 8B        & CoT        & 33.3 & 20.0 & 80.0 & 4.73 \\
\midrule
DeepSeek-Coder 6.7B & Zero-shot  & 20.0 & 20.0 & 80.0 & 4.73 \\
DeepSeek-Coder 6.7B & Few-shot   & 26.7 & 13.3 & 86.7 & 4.60 \\
DeepSeek-Coder 6.7B & CoT        & 40.0 & 26.7 & 73.3 & 4.27 \\
\midrule
Mistral 7B          & Zero-shot  & 40.0 & 73.3 & 26.7 & 2.80 \\
Mistral 7B          & Few-shot   & 26.7 & 86.7 & 13.3 & 2.20 \\
Mistral 7B          & CoT        & 26.7 & 66.7 & 33.3 & 2.87 \\
\bottomrule
\end{tabular}
\label{tab:full_results}
\end{table}

Best condition: \textbf{Mistral 7B + Few-shot} (Bandit@0=26.7\%, Resolution=86.7\%, Avg Iter=2.20). Worst: \textbf{Llama 3.1 8B + Zero-shot} (Bandit@0=66.7\%, Resolution=13.3\%).

\subsection*{CWE-Level Analysis}

\textbf{CWE-78} (OS Command Injection) is the most persistent vulnerability across all conditions, followed by CWE-502 and CWE-434. CWE-89 (SQL Injection) appears in multiple conditions despite being covered by the few-shot examples, indicating model-specific resistance to that fix pattern.

% ─────────────────────────────────────────────────────────────────────────────
\section{Planned Improvements and Technical Direction}
% ─────────────────────────────────────────────────────────────────────────────

\begin{itemize}[noitemsep]
    \item Scale to all 150 LLMSecEval Python prompts for statistical power
    \item Expand few-shot examples to cover CWE-434 and CWE-22, the next most frequent categories
    \item Add a security-aware zero-shot variant (OWASP Top 10 instruction, no examples) to isolate phrasing effects
    \item Measure functional correctness (pass@k on a HumanEval subset) to detect security/functionality trade-off
\end{itemize}

% ─────────────────────────────────────────────────────────────────────────────
\section{Ablation and Prompt Sensitivity Plan}
% ─────────────────────────────────────────────────────────────────────────────

\begin{itemize}[noitemsep]
    \item \textbf{Strategy ablation on Mistral 7B} — isolate whether the resolution advantage is model-intrinsic or strategy-dependent (50 prompts, fixed seed)
    \item \textbf{Few-shot example order} — 3 permutations on 15-prompt subset to measure variance
    \item \textbf{CoT depth} — 2-step vs.\ 4-step CoT on 15-prompt subset
    \item \textbf{Iteration cap} — MAX\_ATTEMPTS=3 vs.\ 5 to assess marginal value of extra repair rounds
\end{itemize}

% ─────────────────────────────────────────────────────────────────────────────
\section{Error Analysis and Generation Quality Analysis}
% ─────────────────────────────────────────────────────────────────────────────

Three recurring failure patterns were identified in files that remained unresolved after 5 iterations:

\begin{enumerate}[noitemsep]
    \item \textbf{Loop divergence} — Bandit passes but the functional checker returns ``Incorrect'' repeatedly (most common in Llama 3.1 8B). The judge appears to apply stricter criteria than the docstring requires (e.g., demanding Flask-WTF instead of \texttt{request.form.get()}).
    \item \textbf{Partial fix} — Model removes one Bandit issue but introduces another (observed in DeepSeek-Coder, CWE-732/CWE-200 co-occurring files).
    \item \textbf{Early exit} — For short functions, some models answer ``Correct'' immediately despite pending SAST warnings, inflating the resolved count.
\end{enumerate}

A full qualitative analysis ($\geq$5 files per model per CWE) will be completed for the final report.

% ─────────────────────────────────────────────────────────────────────────────
\section{Ethical Considerations and Bias}
% ─────────────────────────────────────────────────────────────────────────────

\textbf{Dual-use:} All generated code is stored locally and used only to measure vulnerability rates. LLMSecEval is a published academic benchmark; no novel exploit code is generated.

\textbf{Measurement bias:} Bandit is rule-based and produces both false positives (e.g., \texttt{debug=True} flagged as HIGH in development contexts) and false negatives (logic flaws). The resolution rate therefore measures \emph{Bandit-compliance}, not true security.

\textbf{Model bias:} Mistral 7B's higher resolution rate may partly reflect stronger English instruction-following rather than superior security reasoning. Models differ in training data, instruction-tuning, and language mix.

\textbf{Reproducibility:} Fixed seeds, locally-hosted models (no API non-determinism), and a shared \texttt{data/mid\_phase\_prompts.json} ensure cross-branch comparability.

% ─────────────────────────────────────────────────────────────────────────────
\section{GitHub and Reproducibility Status}
% ─────────────────────────────────────────────────────────────────────────────

Repository: \href{https://github.com/Aliekinozcetin/CodeEnhancer}{github.com/Aliekinozcetin/CodeEnhancer} (\texttt{main} branch). Environment: Python 3.12, \texttt{.venv/}; models via \texttt{ollama pull}; \texttt{.env} gitignored. To reproduce:

\begin{verbatim}
git clone https://github.com/Aliekinozcetin/CodeEnhancer && cd CodeEnhancer
python3 -m venv .venv && source .venv/bin/activate && pip install -r requirements.txt
ollama pull llama3.1 && ollama pull deepseek-coder:6.7b && ollama pull mistral:7b
python3 code_generator.py && python3 code_validator.py
python3 analysis/compare_results.py && python3 analysis/visualize.py
\end{verbatim}

% ─────────────────────────────────────────────────────────────────────────────
\section{Current Challenges and Next Steps}
% ─────────────────────────────────────────────────────────────────────────────

\textbf{Challenges:} (1) Loop divergence in Llama/DeepSeek — same model acting as judge is inconsistent; a dedicated judge model may be needed. (2) Compute time for 150 prompts × 9 conditions × 5 iterations is \textasciitilde{}90–180 hours sequential on M4; batched overnight runs are planned. (3) CWE-200 false positives from logging statements inflate unresolved rates.

\textbf{Next steps:} Scale to all 150 prompts; add per-CWE resolution breakdown in \texttt{compare\_results.py}; qualitative error analysis ($\geq$5 unresolved files per model); draft final report in LNCS template; maintain regular commit history.

% ─────────────────────────────────────────────────────────────────────────────
\begin{thebibliography}{9}

\bibitem{pearce2022asleep}
H.~Pearce, B.~Ahmad, B.~Tan, B.~Dolan-Gavitt, R.~Karri,
``Asleep at the Keyboard? Assessing the Security of GitHub Copilot's Code Contributions,''
\textit{IEEE S\&P}, 2022.

\bibitem{tony2023llmseceval}
C.~Tony, M.~Mutas, N.~Díaz Ferreyra, R.~Scandariato,
``LLMSecEval: A Dataset of Natural Language Prompts for Security Evaluations,''
\textit{MSR}, 2023.

\bibitem{nishida2024codeenhancer}
K.~Nishida, T.~Shibata, H.~Iida,
``CodeEnhancer: Enhancing LLM-Generated Code for Security and Reliability,''
\textit{JAIST Technical Report}, 2024.

\bibitem{brown2020gpt3}
T.~Brown et al.,
``Language Models are Few-Shot Learners,''
\textit{NeurIPS}, 2020.

\bibitem{wei2022cot}
J.~Wei et al.,
``Chain-of-Thought Prompting Elicits Reasoning in Large Language Models,''
\textit{NeurIPS}, 2022.

\end{thebibliography}

\end{document}
